
\section{On the Congregations’ Responsibility for Synodical Acts}

Editorial article in “Folkebladet” for February 24, 1897.

Some time ago there was printed in “Folkebladet” a piece on “The Free-Church Inner Mission.” It pointed out what effect the Ole Nilsen Decision had upon the synod’s mission work, namely, to make it a partisan matter.

On that occasion “Lutheraneren,” in its issue of February 4 of this year, writes:

“No darker accusation can be brought against a synod of congregations. It amounts to nothing less than to pronounce all these congregations devoid of all Christian life. For where there is Christian life, there it always stands also as a principal task to carry the message of salvation in Christ onward to those who do not know it.

If, then, The United Church, which consists of more than 700 congregations, has something else as its principal task, then at any rate the greater part of these 700 congregations must not only be spiritually dead, but so perverted that they do not even care to bear the Christian name.”

From this one may thus see that “Lutheraneren” straightforwardly makes all the congregations of the synod, and all the individual members of them, responsible for the Ole Nilsen Decision to so great a degree that the relation becomes this: if anyone censures this decision and its effect upon The United Church’s missionary work, then that is the same as “pronouncing all the 700 congregations devoid of all Christian life.”

So strong, in the opinion of that organ, is the solidarity in The United Church: if the Ole Nilsen Decision is unchristian, then the synod is unchristian, and then there is no Christian life in 700 congregations?

This claim concerning responsibility goes farther, it seems to us, than is fair. But if within The United Church it is to hold that each congregation and each congregation member is responsible, with person, soul, body, and goods, for all that the annual meeting resolves, then it is a question whether it will be possible for any congregation or person to remain with it, if one otherwise cares what one takes upon oneself.

If the synod adopts a very upright decision in money matters, incurs a vast debt, or the like, will then each congregation and each congregation member be liable, one for all and all for one?

And if the synod adopts an unchristian decision, is that then something which really annihilates the Christian life of all the individual congregation members?

Is there no way out for those who vote against it, testify against it, and work against it either?
Does the Christian life of all the members stand and fall with a single synodical decision?

Yes, so it appears, if the claim of “Lutheraneren” is correct, that he who censures a synodical decision concerning inner mission thereby pronounces 700 congregations devoid of all Christian life.

If the matter of the congregations’ responsibility for the annual meeting’s decisions hangs together in this way, then it is altogether understandable why The United Church excludes delegates and casts out congregations which oppose a proposal or a decision. For if every testimony concerning the annual meeting’s misstep is a complete sentence of condemnation over all congregations and congregation members, then it is altogether natural that such a testimony is not tolerated.

For us this understanding of solidarity is a new and strong testimony that the majority means it sincerely in its declaration that its conduct up to 1895 has been without guilt in the sinful strife in the synod. For only a majority infallible in its own eyes can feel itself so deeply offended by harsh criticism or reproof.

Will the congregations now take this matter seriously and really hold fast to these two things: that, first, they are responsible for all that the annual meeting determines, so that the annual meeting’s decision carries the same moral responsibility for them as if they themselves had determined it;

and that, second, it is sinful and presumptuous for any congregation or any congregation member to oppose or criticize any synodical decision, because by doing so they pronounce 700 congregations devoid of all Christian life?

It is very possible that some congregations would gladly have it so. It is quite convenient to settle down in the assurance that, being so many, one must be right and cannot be mistaken.

But then one must not be surprised either that there still remain some independent-minded congregations and congregation members who will reserve to themselves their own conviction and their own responsibility, and be cast out, if that is not tolerated.

Georg Sverdrup.
